%%%%%%%%%%%%%%%%%%%%%%%%%%%%%%%%%%%%%%%%%%%%%%%%%%%%%%%%%%%%%%%
% START CUSTOM INCLUDES & DEFINITIONS
%%%%%%%%%%%%%%%%%%%%%%%%%%%%%%%%%%%%%%%%%%%%%%%%%%%%%%%%%%%%%%%

%-------------------------------------------------------------
% 1. CORE MATH & FONT PACKAGES
%-------------------------------------------------------------
\usepackage{amsmath,amssymb,amsfonts}
\usepackage{mathtools}
\usepackage{mymath}
\usepackage{MnSymbol}

%-------------------------------------------------------------
% 2. REFERENCE & HYPERLINKING
%-------------------------------------------------------------
\usepackage[hidelinks]{hyperref}
\hypersetup{draft=false}

\usepackage[backend=bibtex, sorting=none, firstinits=true, abbreviate=true, url=false, isbn=false]{biblatex}
\addbibresource{master_bib_abbrev.bib}

%-------------------------------------------------------------
% 3. FIGURES, GRAPHS, & DRAWING (TikZ/PGF)
%-------------------------------------------------------------
\usepackage{tikz}
\usepackage{pgfplots}
\pgfplotsset{compat=newest}
\usetikzlibrary{plotmarks}
\usepgfplotslibrary{groupplots}

\usetikzlibrary{shapes, arrows, shapes.misc, arrows.meta, positioning, matrix, calc, fit, fadings, patterns}
\usetikzlibrary{calc,patterns,decorations.pathmorphing,decorations.markings,arrows, arrows.meta}
\usepackage{varwidth}
\usepackage[export]{adjustbox}
\usepackage{overpic}

%-------------------------------------------------------------
% 4. FLOATING OBJECTS, CAPTIONS, & LAYOUT
%-------------------------------------------------------------
\usepackage{caption}
\usepackage{subcaption}
\usepackage{placeins}

% That's because it likes placing large figures on their own page even though there's space for text
\renewcommand{\floatpagefraction}{0.7}
\renewcommand{\textfraction}{0.1}
\renewcommand{\topfraction}{0.9}
\renewcommand{\bottomfraction}{0.9}

%-------------------------------------------------------------
% 5. TABLES & ALIGNMENT
%-------------------------------------------------------------
\usepackage{arydshln}
\usepackage{tabto}
\usepackage{siunitx}
\usepackage{nth}

%-------------------------------------------------------------
% 6. ALGORITHMS & PSEUDOCODE
%-------------------------------------------------------------
\usepackage{algorithm}
\usepackage[noend]{algpseudocode}
\usepackage{xcolor}

\definecolor{commentgrey}{gray}{0.6}
\newcommand{\algcomment}[1]{\scriptsize\textcolor{commentgrey}{\Comment{#1}}}

%-------------------------------------------------------------
% 7. THEOREM-LIKE ENVIRONMENTS
%-------------------------------------------------------------
\newtheorem{assumption}{Assumption}
\newtheorem{remark}{Remark}

%-------------------------------------------------------------
% 8. UTILITY & CLASS FILE MODIFICATIONS
%-------------------------------------------------------------
\usepackage{todonotes}

% Remove preprint footer
\makeatletter
\def\ps@copyright{\let\@mkboth\@gobbletwo
\def\@oddhead{}
\let\@evenhead\@oddhead
\def\@oddfoot{\hfil\thepage\hfil}
\let\@evenfoot\@oddfoot}
\makeatother

% Patch to close NoHyper scope
\edef\endfrontmatter{%

\unexpanded\expandafter{\endfrontmatter}% the current code
\noexpand\endNoHyper
}

%-------------------------------------------------------------
% 9. CUSTOM MATH OPERATORS & COMMANDS
%-------------------------------------------------------------
\newcommand{\Argmin}{\ensuremath{\mathop{\mathrm{argmin}}}}
\newcommand{\Argmax}{\ensuremath{\mathop{\mathrm{argmax}}}}
\newcommand{\twonorm}[1]{\left\| #1 \right\|_2}
\newcommand{\set}[2]{\left\{ #1 \;\middle|\; #2 \right\}}
\newcommand{\inR}[3]{#1 \in \Rset^{#2 \times #3}}
\newcommand{\trans}[1]{#1^\top}

%%%%%%%%%%%%%%%%%%%%%%%%%%%%%%%%%%%%%%%%%%%%%%%%%%%%%%%%%%%%%%%
% END CUSTOM INCLUDES & DEFINITIONS
%%%%%%%%%%%%%%%%%%%%%%%%%%%%%%%%%%%%%%%%%%%%%%%%%%%%%%%%%%%%%%%